\documentclass[10pt,a4paper]{amsart}
\usepackage[margin=19mm]{geometry}
\usepackage{amsmath,amssymb,mathtools}
\usepackage[scr=rsfs,cal=euler]{mathalfa}
\usepackage{xfrac}
\usepackage[unicode,hidelinks,bookmarks=false]{hyperref}
\newcommand{\ie}{i.e.\ }
\newcommand{\cf}{cf.\ }
\newcommand{\resp}{respectively\ }
\newcommand{\Subsection}{\subsection}
\providecommand{\nobreakdash}{\nobreak\hbox{-}}
\setlength{\emergencystretch}{2em}
\multlinegap0pt
\usepackage{bm}
\usepackage{bbm}
\usepackage{dsfont}
\usepackage{xspace}
\usepackage{cancel}
\usepackage{mathtools}
\usepackage{amsthm}
\makeatletter
\@ifundefined{thm}{\newtheorem{thm}{Theorem}}{}
\makeatother
\newcommand{\cl}{\mbox{{\rm cl}}}
\newcommand{\pr}{\mbox{{\rm pr}}}
\title[On two methods of constructing compactifications of topologi]{On two methods of constructing compactifications of topological groups}
\author{K. L. Kozlov, A. G. Leiderman}
\date{Source: arXiv:2507.10353v2 (CC BY 4.0)}
\begin{document}
\maketitle

\noindent\emph{Source and licence.} On two methods of constructing compactifications of topological groups. Version arXiv:2507.10353v2 (CC BY 4.0), distributed under the Creative Commons Attribution 4.0 International licence, as recorded in the arXiv licence metadata for this exact version and shown on the abstract page https://arxiv.org/abs/2507.10353v2. Full rights notice: licenses/arxiv-2507.10353-CC-BY-4.0.txt.\par\smallskip

\noindent\textit{Editorial scope.} Counted unit: the Thm whose source label is \texttt{thm1} in the article (source0.tex lines 843--850, source label \texttt{thm1}), delivered together with the article's own complete original proof of it (source0.tex lines 852--893, one \texttt{proof} environment of 42 lines). No mathematics is written, repaired or re-derived: statement and proof are the article's own text line for line. Required context retained uncounted: none. Every definition, display and label of those lines is retained unchanged. Omitted from this package and not redistributed: 1--842, 851--851, 894--1214. Nothing inside a retained excerpt is omitted or altered: every retained body is delivered exactly as the article's lines are written, so no omission receipt of any kind accompanies this package. Citations: the retained excerpts carry no citation command at all, so no bibliography entry of the article is transcribed and no reference is redistributed. Reference adaptation: none was needed and none was made; every reference inside the delivered unit resolves inside it, so no unresolved reference or citation is produced. Printed slips kept literal: no printed word, symbol, hypothesis or number of the counted statement or of its proof was altered. Layout and class: the article's own class is \texttt{amsart} with packages including amsmath, latexsym, amscd, amsbsy, amssymb, amsfonts, amsthm, fleqn, leqno, euscript, graphicx, color; the wrapper is the lane's pinned \texttt{amsart} editorial wrapper at 10pt on A4, which restates the theorem-like environments the retained text needs and adds the article's own macro definitions that the retained text needs (\texttt{\textbackslash cl}, \texttt{\textbackslash pr}), with extra wrapper packages (bm, bbm, dsfont, xspace, cancel, mathtools). The delivered numbering is the unit's own. Assets: no retained span contains a figure environment, a graphics-inclusion command, a PDF-page inclusion, a table or a TikZ picture; no image file is copied into this package, the package declares no asset dependency and no image file is redistributed with it. Licence: version arXiv:2507.10353v2 of this article is distributed under the Creative Commons Attribution 4.0 International licence, as recorded in the arXiv licence metadata for this exact version and shown on the abstract page https://arxiv.org/abs/2507.10353v2. A full rights notice is delivered at licenses/arxiv-2507.10353-CC-BY-4.0.txt. Characters: every retained excerpt is pure ASCII, so the delivered engine, XeLaTeX with its default Latin Modern text font, reports no missing character.
\begin{thm}\label{thm1}
The Ellis map $\mathfrak{E}:\mathbb{G}(G)\to \mathbb{E} (G)$ has the following properties
\begin{itemize}
\item $\mathfrak{E}$ is an epimorphism {\rm(}of posets{\rm)}, 
\item $\forall\ b G\in \mathbb{G} (G)$  $\mathfrak{E}(b G)\geq b G$, 
\item $\mathfrak{E}\circ\mathfrak{E}=\mathfrak{E}$.
\end{itemize}
\end{thm}

\begin{proof}
(i)  $\mathfrak{E}$ is morphism.  Let $(b' G, b')$, $(b G, b)$ be $G$-compactifications of $G$, $\varphi: b' G\to b G$ is the $G$-map of $G$-compactifications ($\tilde\alpha$ and $\tilde\alpha'$ are actions of $G$ on $b G$ and $b' G$ respectively). The $G$-map $\varphi$ induces the $G$-map  $\Phi:  (b' G)^{b' G}\to  (b G)^{b G}$ and the following diagram 
$$\begin{array}{ccl}
(b' G)^{b' G} & \stackrel{\Phi}{\longrightarrow} &  (b G)^{b G} \\
\imath_{b' G} \nwarrow &   &  \nearrow \imath_{b G} \\
 & G & 
\end{array}$$
is commutative.  In fact, for any  $h\in G$ 
$$\imath_{b G} (g)(b(h))=\tilde\alpha (g, b(h))=b(\alpha (g, h))\ \mbox{and}$$ 
$$[(\Phi\circ\imath_{b' G}) (g)] (b'(h))=\pr_{b'(h)}\big((\Phi\circ\imath_{b' G}) (g)\big)=\pr_{b'(h)}\big((\varphi_{Y}\circ\varphi^{b' G}\circ \imath_{b' G})(g)\big)=$$
\centerline{($\pr_{b'(h)}\circ\varphi_{Y}=\pr_{b(h)}$)}
$$=\pr_{b(h)}\big((\varphi^{b' G}\circ \imath_{b' G})(g)\big)=\varphi\big(\imath_{b' G} (g)(b'(h))\big)=\varphi\big(\tilde\alpha'(g, b'(h))\big)=$$
\centerline{($\varphi$ is a $G$-map)} 
$$=\tilde\alpha(g, \varphi (b'(h)))=\tilde\alpha(g, b(h))\stackrel{h\in G}=b(\alpha(g, h)).$$
Thus, $\imath_{b G} (g)(t)=(\Phi\circ\imath_{b' G}) (g)(t)$, $t\in b (G)$.  Being continuous they coincide on $b G$ and  $\imath_{b G} (g)=(\Phi\circ\imath_{b' G}) (g)$, $g\in G$.

Continuity of $\Phi$, commutativity of the above diagram ($\Phi|_{\imath_{b' G}}$ is a homeomorphism) and compactness of $\cl~\imath_{b' G} (G)$, $\cl~\imath_{b G} (G)$ yields that the restriction $\Phi|_{\mathfrak{E}(b' G)}: \mathfrak{E}(b' G)\to \mathfrak{E}(b G)$ is a map of $G$-compactifications. Hence, $\mathfrak{E}(b' G)\geq\mathfrak{E}(b G)$ and $\mathfrak{E}$ is a morphism of posets. 

\medskip

(ii) Let $(b G, b)$ be $G$-compactification of $G$. The projection $\pr_{e}: b G^{b G}\to b G=\{e\}\times b G$ is continuous and 
$$\pr_{b(e)}\circ\imath_{b G}=b.$$
Indeed, $\pr_{b(e)}\circ\imath_{b G}(g)=\imath_{b G}(g)(b(e))=\tilde\alpha (g, b(e))=b(g)$, $g\in G$. The restriction of $\pr_{b(e)}$ to $\imath_{b G} (G)$ is a homeomorphism. Hence,  $\pr_{b(e)}|_{\mathfrak{E}(b G)}$ is the map of compactifications and $\mathfrak{E}(b G)\geq b G$. 

\medskip

(iii) By the previous $\mathfrak{E}(\mathfrak{E}(b G))\geq\mathfrak{E} (b G)$ and   $\mathfrak{E}(b G)\in\mathbb{E} (G)$. It remains to prove that if $(b G, b)$ is an Ellis compactification of $G$, then $b G\geq\mathfrak{E}(b G)$. 

Let $\tilde\alpha: G\times bG\to bG$ be the extension of action $\alpha$, $\imath_{b G}: G\to (bG)^{bG}$ is an embedding. 

Since $(bG, \bullet)$ is a right topological monoid and continuous multiplication on the right is uniformly continuous on the compact space $b G$ (with the unique uniformity $\tilde{\mathcal U}$), from 
$$\imath_{b G} (g)(t)=\tilde\alpha (g, t)=b (g)\bullet t,\ t\in bG,\ g\in G,$$
it follows that for any $t\in bG$ and   for any ${\rm U}\in \tilde{\mathcal U}$ there exists ${\rm V}\in \tilde{\mathcal U}$ such that 
$$\mbox{if for}\ f, g\in G,\ (f, g)\in {\rm V},\ \mbox{then}\ (b (f)\bullet t, b (g)\bullet t)\in {\rm U}.$$
Therefore, for any entourage $\hat {\rm U}=\{ (f, g)\in G\times G\ |\ (\imath_{b G} (f)(t)=b (f)\bullet t, \imath_{b G} (g)(t)=b (g)\bullet t)\in {\rm U}\}$, where $t\in bG$, ${\rm U}\in\tilde{\mathcal U}$ (from the subbase of the subspace uniformty on $\imath_{b G}(G)$ as a subspace of $b G^{b G}$), there exists  ${\rm V}\in\tilde{\mathcal U}$ such that 
$$\mbox{if}\ (f, g)\in {\rm V},\ \mbox{then}\  (f, g)\in \hat {\rm U}$$
and the map $\imath_{b G}: G\to\imath_{b G} (G)$ is uniformly continuous. Therefore,  $\imath_{b G}$ can be extended to the map of compactifications $\tilde\imath_{b G}: b G\to\mathfrak{E}(b G)$. Hence, $b G\geq\mathfrak{E}(b G)$ and $\mathfrak{E}\circ\mathfrak{E}=\mathfrak{E}$.

\medskip

(i) $\mathfrak{E}$ is an epimorphism. If $b G\in\mathbb{E}(G)$, then $b G=\mathfrak{E}(b G)$. Indeed, by (i) $b G\leq\mathfrak{E}(b G)$, by (ii) $b G\geq\mathfrak{E}(b G)$. Hence, $b G=\mathfrak{E}(b G)$. From this it follows that $\mathfrak{E}$ is an epimorphism. 
\end{proof}

\end{document}
